\documentclass[11pt,a4paper]{article}
\usepackage[margin=2.5cm]{geometry}
\usepackage{amsmath,amssymb,booktabs}
\usepackage[hidelinks]{hyperref}
\title{Whole genome search for nested tandem repeats: work log}
\author{Atheer Matroud}
\date{Sessions of 11 November 2026}
\begin{document}
\maketitle

\section{Aim}
The aim is a new software tool that searches a whole genome for nested tandem
repeats (NTRs) efficiently. NTRFinder (Chapter 4 of the thesis) first finds a
tandem repeat with two scanning windows and then looks for displaced copies of
the tandem motif. The new tool replaces that search phase by statistics on the
distances between occurrences of $k$-tuples, collected in one pass over the
sequence. Verification remains the nested wrap-around dynamic programming of
Chapter 5.

\section{Decisions taken in this session}
\begin{enumerate}
\item The input is a FASTA file that may hold a whole genome.
\item The user chooses a motif length interval. Three intervals are provided,
      each with its own parameters: 7 to 100 bp, 101 to 500 bp, and 501 bp and
      above.
\item The tuple length grows with the interval: $k=4$, $k=6$, $k=8$.
\item A locus may contain more than two motifs. Verification will use the
      $r$-motif alignment of Section 5.4.4 of the thesis.
\item The implementation is in Python with NumPy and Numba.
\item The sequence is scanned from left to right and every overlapping
      $k$-tuple is visited.
\item The work proceeds step by step, with tests on simulated and real data at
      each step.
\end{enumerate}

\section{Planned steps}
\begin{enumerate}
\item Read the FASTA file record by record and encode it.
\item Compute, for every site, the distance to the previous occurrence of the
      $k$-tuple that starts there.
\item Compress the distances into runs of constant value and discard short runs.
\item Chain neighbouring runs into loci and classify each locus as a tandem
      repeat or a candidate NTR.
\item Predict the motifs: the tandem motif $\mathbf{x}$ from the dominant
      distance, the interspersed motifs by clustering the segments between runs
      of $\mathbf{x}$.
\item Verify by $r$-motif nested wrap-around dynamic programming.
\item Report the loci.
\item Test on simulated data (Section 4.4.1 of the thesis) and on real data.
\end{enumerate}
Steps 1 and 2 were completed in this session.

\section{The distances in a nested tandem repeat}
Let $\mathbf{x}$ be the tandem motif with $|\mathbf{x}|=m$ and $\mathbf{X}$ the
interspersed motif with $|\mathbf{X}|=l$, and let the NTR be
$\mathbf{x}^{s_0}\mathbf{X}\mathbf{x}^{s_1}\mathbf{X}\cdots\mathbf{X}\mathbf{x}^{s_n}$.
For a $k$-tuple that starts at site $i$, let $d(i)$ be the distance from $i$ back
to the start of the previous occurrence of the same $k$-tuple. Then
\begin{center}
\begin{tabular}{ll}
\toprule
The $k$-tuple lies in & $d(i)$ \\
\midrule
$\mathbf{x}$, inside a run $\mathbf{x}^{s_i}$ & $m$ \\
the first $\mathbf{x}$ after an $\mathbf{X}$ & $m+l$ \\
$\mathbf{X}$ & $l+s_i m$ \\
\bottomrule
\end{tabular}
\end{center}
Every distance other than $m$ is congruent to $l$ modulo $m$. The dominant
distance gives $m$, the smallest other distance gives $m+l$, and the
congruence is a test that costs no alignment. With $r$ motifs of lengths
$l_1,\ldots,l_r$ every distance has the form $\sum_j a_j l_j$ with non-negative
integers $a_j$.

\section{Step 1: reading the FASTA file}
The file is read one record at a time, so only one sequence is held in memory.
The nucleotides are encoded as $A=0$, $C=1$, $G=2$, $T=3$. Every other symbol is
marked invalid. Lower case letters are treated as upper case. Compressed
(\texttt{.gz}) files are accepted. Code: \texttt{ntrscan/fasta.py}.

\section{Step 2: the distances $d(i)$}
The sequence is scanned once from left to right. The $k$-tuple that starts at
site $i$ covers sites $i,\ldots,i+k-1$, and the next $k$-tuple starts at site
$i+1$, so consecutive tuples overlap in $k-1$ sites and none is skipped. A table
with $4^k$ entries holds the start of the last occurrence of every $k$-tuple.
For each $i$ the table gives $d(i)$, and the entry is then replaced by $i$.
$d(i)=0$ is recorded when there is no previous occurrence within the largest
distance allowed for the interval, or when the $k$-tuple contains an invalid
symbol. The scan takes $O(n)$ time and $O(4^k)$ additional space. Code:
\texttt{ntrscan/track.py}.

The provisional parameters are given in Table~\ref{tab:tiers}.
\begin{table}[h]\centering
\caption{Provisional parameters of the three intervals.}\label{tab:tiers}
\begin{tabular}{lrrr}
\toprule
Motif length (bp) & $k$ & Largest distance recorded \\
\midrule
7 to 100 & 4 & 2000 \\
101 to 500 & 6 & 10000 \\
501 and above & 8 & 65535 \\
\bottomrule
\end{tabular}
\end{table}

\subsection{Worked example}
With $k=4$, $\mathbf{x}=\texttt{ACGGTCA}$ ($m=7$),
$\mathbf{X}=\texttt{GACCTTTGGC}$ ($l=10$), the NTR
$\mathbf{x}^3\mathbf{X}\mathbf{x}^2\mathbf{X}\mathbf{x}^4\mathbf{X}\mathbf{x}^2$
and 12 random nucleotides on each side, the runs of constant $d$ are given in
Table~\ref{tab:example}. Sites are numbered from 0 and the NTR starts at site 12.
\begin{table}[h]\centering
\caption{Runs of constant $d$ in the worked example (short runs omitted).}\label{tab:example}
\begin{tabular}{rrrl}
\toprule
Sites & Run length & $d$ & Meaning \\
\midrule
18 to 29 & 12 & 7 & inside $\mathbf{x}^3$ \\
43 to 46 & 4 & 17 & first $\mathbf{x}$ after $\mathbf{X}$: $m+l$ \\
47 to 49 & 3 & 24 & tuples that span two copies of $\mathbf{x}$ \\
50 to 53 & 4 & 7 & inside $\mathbf{x}^2$ \\
54 to 66 & 13 & 24 & second $\mathbf{X}$: $l+2m$ \\
67 to 70 & 4 & 17 & $m+l$ \\
74 to 91 & 18 & 7 & inside $\mathbf{x}^4$ \\
92 to 104 & 13 & 38 & third $\mathbf{X}$: $l+4m$ \\
105 to 108 & 4 & 17 & $m+l$ \\
112 to 116 & 5 & 7 & inside $\mathbf{x}^2$ \\
\bottomrule
\end{tabular}
\end{table}
The dominant distance is 7, so $m=7$. The smallest other distance is 17, so
$l=10$. The distances 24 and 38 leave remainder 3 on division by 7, as $l$
does. The runs of length $13=l+k-1$ mark the copies of $\mathbf{X}$: the run at
sites 54 to 66 places $\mathbf{X}$ at sites 57 to 66, where it was embedded.
The multipliers 2 and 4 in $l+2m$ and $l+4m$ are the values $s_i$.

With a third motif $\mathbf{Y}=\texttt{TTAGCATCGAAC}$ (12 bp) and the order
$\mathbf{x}^3\mathbf{X}\mathbf{Y}\mathbf{x}^2\mathbf{X}\mathbf{x}^3\mathbf{Y}^2\mathbf{x}^2\mathbf{X}\mathbf{x}^2$
the distances observed are 7, 12, 17, 24, 29, 31, 36, 38, 57 and 69, each of the
form $7a+10b+12c$.

\section{Tests of Steps 1 and 2}
\subsection{Correctness}
Seven tests are in \texttt{tests/test\_step1\_step2.py} and all pass:
encoding; reading plain and compressed FASTA files with several records;
overlapping tuples on hand-checked strings; invalid symbols, lower case and the
largest distance; agreement with a slow dictionary implementation on 200 random
sequences for $k\in\{1,2,3,4,6,8\}$; the worked example of
Table~\ref{tab:example}; and 100 exact NTRs embedded in 100\,000 random
nucleotides.

\subsection{Simulated data}
The design of Section 4.4.1 of the thesis was followed: random sequence with
each nucleotide at probability 0.25, 100 embedded NTRs with random motifs of 7
to 100 bp, indels at 1\% per site and substitutions at 1\% to 4\% per site.
Table~\ref{tab:signal} gives the fraction of the tandem sites of $\mathbf{x}$
at which $d(i)=m$ (exact data) or $|d(i)-m|\le 1$ (mutated data).
\begin{table}[h]\centering
\caption{Fraction of tandem sites of $\mathbf{x}$ with $d(i)=m$, by motif length.}\label{tab:signal}
\begin{tabular}{rrrrrrr}
\toprule
$k$ & Substitutions & 7 to 20 & 21 to 40 & 41 to 60 & 61 to 80 & 81 to 100 \\
\midrule
4 & 0\% & 0.85 & 0.87 & 0.85 & 0.73 & 0.68 \\
4 & 1\% & 0.75 & 0.78 & 0.73 & 0.58 & 0.55 \\
4 & 2\% & 0.69 & 0.72 & 0.67 & 0.55 & 0.50 \\
4 & 3\% & 0.63 & 0.66 & 0.60 & 0.51 & 0.47 \\
4 & 4\% & 0.65 & 0.60 & 0.58 & 0.47 & 0.45 \\
\midrule
6 & 0\% & 0.87 & 0.94 & 0.97 & 0.95 & 0.96 \\
6 & 1\% & 0.70 & 0.77 & 0.76 & 0.71 & 0.71 \\
6 & 2\% & 0.61 & 0.68 & 0.67 & 0.64 & 0.62 \\
6 & 3\% & 0.52 & 0.61 & 0.58 & 0.58 & 0.56 \\
6 & 4\% & 0.54 & 0.52 & 0.53 & 0.50 & 0.50 \\
\midrule
8 & 0\% & 0.83 & 0.92 & 0.95 & 0.96 & 0.97 \\
8 & 1\% & 0.61 & 0.69 & 0.69 & 0.66 & 0.67 \\
8 & 2\% & 0.51 & 0.60 & 0.58 & 0.57 & 0.56 \\
8 & 3\% & 0.41 & 0.51 & 0.49 & 0.50 & 0.48 \\
8 & 4\% & 0.43 & 0.41 & 0.43 & 0.41 & 0.42 \\
\bottomrule
\end{tabular}

\smallskip
{\footnotesize The rows with substitutions also carry indels at 1\% per site.}
\end{table}

Two observations follow.
\begin{enumerate}
\item In exact repeats the fraction is below 1 for two reasons. Tuples that
      span the end of the last $\mathbf{x}$ of a run have a different distance.
      A $k$-tuple that occurs twice inside $\mathbf{x}$ gives a distance
      shorter than $m$. The second effect grows with $m$: with $k=4$ the
      fraction falls to 0.68 for motifs of 81 to 100 bp, while $k=6$ stays at
      0.96.
\item A larger $k$ loses more sites to mutation, because one substitution
      destroys $k$ tuples. At 4\% substitutions about half of the sites keep
      the signal for $k=4$ and $k=6$.
\end{enumerate}

In 20 Mb of random sequence, runs of constant non-zero $d$ of length 10 or more
occurred 10, 7 and 12 times for $k=4,6,8$, and the longest run was 10, 11 and 12
sites. A minimum run length near 12 is therefore a candidate threshold for Step
3, to be confirmed on real genomes.

\subsection{Speed}
On 50 Mb the scan alone took 0.07 s, 0.11 s and 0.26 s for $k=4,6,8$. Reading a
50 Mb FASTA file and scanning it with $k=4$ took 0.3 s in total, which
corresponds to about 21 s for a genome of 3.1 Gb on two cores of a cloud
machine.

\subsection{Real data}
Real data were tested after the simulated data, in the manual investigation of
Section~\ref{sec:manual}.

\section{Further rulings}
\begin{enumerate}
\item The interval 7 to 100 bp is split: 7 to 60 bp with $k=4$ and 61 to 100 bp
      with $k=6$. Section~\ref{sec:taro} shows that this ruling needs to be
      revisited.
\item A locus in which two different halves alternate strictly, so that every
      $s_i=1$, is reported as a tandem repeat with motif $\mathbf{xX}$.
\item A locus is an NTR only if $d(\mathbf{x},\mathbf{X})$ exceeds a threshold
      $h$. The threshold is a user parameter; the default is equation (4.1) of
      the thesis with $p_m=0.8$.
\item An NTR with a short tandem motif inside a long interspersed motif is
      reported only if $|\mathbf{x}|>1$ and the $s_i$ are not all equal.
\item The results are examined by hand first. The pipeline is built once the
      observations are correct.
\end{enumerate}

\section{Manual investigation on real data}\label{sec:manual}
\subsection{Human Y chromosome}
The file \texttt{HC18989.HIFIRW.ONTUL.na.chrY.fasta} holds 44 records and about
71 Mb, with no \texttt{N}. The scan takes about one second per 10 Mb. For the
largest record (PAR1 to centromere, 10.6 Mb) and $k=4$, runs of constant
non-zero $d(i)$ of length 12 or more cover 0.88\% of the record (4\,538 runs);
runs of 100 or more number 16. The distances that cover most sites are 4, 1, 2,
6, 5 and 3; above 7 bp they are 26, 27, 61, 45, 55, 49 and 41. The 23
centromeric records give no run of 10 or more with $k=4$.

Runs of 12 or more sites with $d\ge 7$ were grouped into loci when they lay
within 200 bp. Of 1\,464 loci, 1\,371 show one distance, 37 show distances that
differ by 1 or 2, 3 show a distance and its multiple, and 53 show mixed
distances. The 53 were read by hand and fall into four groups
(Table~\ref{tab:ygroups}). About 30 of them are shorter than 200 bp and cannot
be judged.

\begin{table}[h]\centering
\caption{Mixed loci on the Y chromosome ($k=4$), record 0.}\label{tab:ygroups}
\begin{tabular}{p{4.2cm}p{5.6cm}p{4.2cm}}
\toprule
Group & Example & Reading \\
\midrule
One motif, copies of different length & site 1\,755\,586: 33 and 28; the short copy lacks \texttt{TAGCC} & tandem repeat \\
Two halves in strict alternation & site 2\,020\,824: 32, and 29 where \texttt{GCC} is deleted & tandem repeat, motif $\mathbf{xX}$ \\
Diverged copies & site 314\,131: 40, and 120 = 3 $\times$ 40 & tandem repeat \\
Short repeat inside a long unit & site 147\,490: 113, 131, 121, each followed by a run with $d=2$ & NTR with $\mathbf{x}=\texttt{TG}$ \\
\bottomrule
\end{tabular}
\end{table}

The locus at sites 1\,949\,408 to 1\,949\,946 shows the distances 21 and 15 and
matches the last row of Table 4.2 of the thesis ($|\mathbf{x}|=21$,
$|\mathbf{X}|=15$). Its copies are
$\mathbf{x}=\texttt{GGACCCAGGCAGATAGGAATA}$ and a 15 bp copy that is
$\mathbf{x}$ without \texttt{GGAATA}. Whether this is an NTR under the
threshold ruling is not yet decided.

With $k=6$ and $d\ge 61$ the centromeric records are dominated by $d=171$
(142\,085 sites), record 24 holds an array of about 250 kb with $d=125$,
record 0 holds an array of about 44 kb with $d=61$, and the \texttt{TG} locus
appears as one locus of 4.2 kb with distances from 110 to 133. Half of the loci
(1\,637 of 3\,323) show two or more distances, mostly $m$, $m\pm 1$ and
multiples of $m$. The mixture of distances alone therefore does not identify an
NTR; the motifs have to be compared.

\subsection{Taro}\label{sec:taro}
The data are two rDNA regions of cultivar Lipu (CM132705.1, sites 131\,800\,001
to 131\,825\,000, and CM132707.1, sites 1 to 25\,200) and ten extracted NTR
alleles from cultivars Longxiangyu, Moi and Lipu. Run tables for $k=4$ and
$k=6$ were supplied with the regions; the scan reproduces all four tables
exactly (23, 71, 16 and 60 runs).

The motifs, read from the sequence, are as in the thesis, with $m=11$ and
$l=48$:
\begin{align*}
\mathbf{x}&=\texttt{GTCGCACAGCC},\\
\mathbf{X}&=\texttt{GTTCTGGGCAAAACGGCTGGGCGACGTGCTGGACTGGCCAGCTGGTTC}.
\end{align*}

\paragraph{The effect of $k$.} Table~\ref{tab:tarok} gives the distances seen
on the first region. With $k=4$ only $d=11$ is seen: the tandem motif is found
and the interspersed motif is not. The distances $l+s_im$ appear from $k=5$ and
are complete at $k=6$. The interspersed motif has 48 bp, which the current
ruling assigns to $k=4$. The relevant length for the choice of $k$ is the
distance $l+s_im$ (59 to 213 here), not $|\mathbf{X}|$.

\begin{table}[h]\centering
\caption{Distances in runs of 12 or more sites, first taro region ($d\ge 7$).}\label{tab:tarok}
\begin{tabular}{rp{11cm}}
\toprule
$k$ & Distances \\
\midrule
4 & 11, 29 \\
5 & 11, 29, 41, 59, 167, 180, 213 \\
6 & 11, 22, 29, 59, 81, 110, 114, 125, 147, 158, 167, 180, 194, 213 \\
8 & 11, 22, 29, 59, 81, 110, 114, 125, 136, 147, 158, 167, 180, 213, 228, 338, 360, 371, 676, 687 \\
\bottomrule
\end{tabular}
\end{table}

\paragraph{Structure.} The copies of $\mathbf{X}$ were located in each allele
by approximate matching and the $s_i$ were read from the gaps between them.
For allele A07 (Lipu) this gives 12 copies of $\mathbf{X}$ and
\[(s_0,\ldots,s_{12})=(4,3,1,6,7,8,4,10,12,10,15,13,4),\]
with $\sum s_i=97$, in agreement with the counts in the header of the allele
($X=12$, $x=97$). The predicted distances are 11, $m+l=59$, and
$l+s_im\in\{81,92,114,125,136,158,180,191,213\}$. The scan with $k=6$ reports
11, 59, 81, 114, 125, 158, 180, 194 and 213. The values 92 and 136 are missing
because the copies of $\mathbf{X}$ concerned carry 3 to 5 differences and their
runs are shorter than 12; 194 replaces 191 because of an indel. Every distance
other than 11 leaves remainder 4 on division by 11, as $l=48$ does. The other
alleles behave in the same way (Table~\ref{tab:taroalleles}). The distance 167
belongs to a different repeat downstream of the NTR.

\begin{table}[h]\centering
\caption{Structure of the taro alleles read from the sequence.}\label{tab:taroalleles}
\begin{tabular}{llrl}
\toprule
Allele & Cultivar & Copies of $\mathbf{X}$ & $(s_0,\ldots,s_n)$ \\
\midrule
A01, A02, A07, A09 & Longxiangyu, Lipu & 12 & (4,3,1,6,7,8,4,10,12,10,15,13,4) \\
A03 & Longxiangyu & 12 & (4,3,1,6,7,8,5,9,12,10,15,13,4) \\
A08 & Lipu & 12 & (4,3,1,6,7,8,4,10,12,9,15,13,4) \\
A04, A06 & Longxiangyu, Moi & 12 & (8,3,2,8,10,9,4,11,6,10,10,6,5) \\
A05 & Moi & 11 & (5,3,1,6,10,5,10,8,10,13,12,4) \\
\bottomrule
\end{tabular}

\smallskip
{\footnotesize The $s_i$ are gap lengths divided by 11 and rounded; they are a
reading by hand, not an alignment. Allele A10 is incomplete in the supplied
file and the file of allele maps is empty.}
\end{table}

\paragraph{Distance between the motifs.} The edit distance
$d(\mathbf{x},\mathbf{X})$ is 37, but 37 is also the difference in length, so
this value says little. The best match of $\mathbf{x}$ against a substring of
$\mathbf{X}$ has 4 differences, against a bound of 3.5 from equation (4.1) for
$|\mathbf{x}|=11$. For the Y chromosome locus with lengths 21 and 15 the short
motif matches a substring of the long one with no difference. A comparison of
the shorter motif with the best substring of the longer one therefore separates
the two cases, whereas the plain edit distance does not.

\section{Open points}
\begin{enumerate}
\item The choice of $k$ should follow the distances $l+s_im$ and not the motif
      length; the split of the first interval has to be reconsidered.
\item The measure used for $d(\mathbf{x},\mathbf{X})$ when the motifs differ in
      length, and the status of the 21 and 15 locus.
\item The minimum run length and the minimum size of a locus. On real data only
      15\% to 40\% of the sites of a locus lie in runs of 12 or more.
\item The largest distance recorded for each interval is provisional.
\end{enumerate}

\section{Note on the code (18 November 2026)}
The package files written on 11 November 2026 (\texttt{ntrscan/fasta.py},
\texttt{ntrscan/track.py} and \texttt{tests/test\_step1\_step2.py}) were not kept.
They were rebuilt on 18 November 2026 from this log and from
\texttt{Taro\_NTR/scan\_taro.py}. The seven tests pass, and the rebuilt scan
reproduces the four taro run tables exactly (23, 71, 16 and 60 runs, with the
largest distance 2000 and a minimum run length of 12). The figures of
Table~\ref{tab:signal} and of the speed test were obtained with the original
files and were not rerun.
\end{document}
